\chapter{\IfLanguageName{dutch}{\gls{PoC}}{\gls{PoC}}}%
\label{ch:proof-of-concept}

In dit hoofdstuk wordt de praktische uitwerking van de bachelorproef gepresenteerd als een \gls{PoC}. 
Waar de voorgaande hoofdstukken dit onderwerp theoretisch kaderen, richt dit hoofdstuk zich op de technische implementatie.

Deze \gls{PoC} dient als een prototype dat de technische haalbaarheid van een geautomatiseerd, schaalbaar notificatiesysteem voor de Vlaamse concertmarkt bewijst.
Hierbij staat de generieke aanpak centraal, de scraping-logica moet toepasbaar zijn voor de verschillende HTML-structuren van de venue pagina's zodat de notificaties vanuit één punt beheerbaar zijn.

\section{Architectuuroverzicht}

Om tot een robuust en schaalbaar systeem te komen, is de server architectuur modulair opgebouwd, waarbij de verschillende processen zoals data-extractie, \textit{change detection}, notificaties en API-implementatie onafhankelijk opereren.
Deze aanpak is cruciaal om de detectiesnelheid, betrouwbaarheid en schaalbaarheid van het systeem te waarborgen.

De modulariteit van de server is duidelijk afgebakend door de mappenstructuur:

\begin{itemize}
    \item \textbf{app/api}: Bevat de REST-endpoints (FastAPI) voor het beheren van concerten, venues, performers en gebruikersvoorkeuren. Dit is de interface waarmee de frontend communiceert.
    \item \textbf{app/core}: Bevat de kernconfiguratie van de applicatie, inclusief de databaseverbinding (SQLAlchemy), globale instellingen en de configuratie voor Celery zodat de achtergrondtaken correct worden uitgevoerd.
    \item \textbf{app/models}: Definieert de database-tabellen en hun onderlinge relaties.
    \item \textbf{app/schemas}: Bevat de Pydantic-modellen voor datavalidatie en conversie (\gls{JSON}) bij API-verzoeken.
    \item \textbf{app/services}: Bevat de business logica van de applicatie, zoals de \textit{matcher} (die concerten koppelt aan gebruikersvoorkeuren) en de \textit{mailer} voor het versturen van e-mailnotificaties.
    \item \textbf{scraper/}: Bevat de gespecialiseerde scraping-engine. Het maakt gebruik van een generiek model om data van verschillende venue-websites te extraheren, te valideren en klaar te zetten voor verwerking.
    \item \textbf{tests/}: Bevat testen waarmee de functionaliteit van de applicatie wordt gewaarborgd en regressies worden voorkomen.
\end{itemize}

\section{Datamodellen en Validatie}

De datamodellen dienen als een uniform bruikbaar formaat, zodat alle verscheidene HTML-structuren consistent kunnen worden verwerkt.
De transformatie van ongestructureerde HTML-data naar gestructureerde datamodellen is een cruciale stap in het proces en bestaat uit volgende fasen:

\subsection{Database Modellen (SQLAlchemy)}

De database architectuur is ontworpen om de relationele verbanden tussen concerten, locaties (\textit{Venues}) en artiesten (\textit{Performers}) te bewaken.
Met behulp van SQLAlchemy worden deze entiteiten omgezet in Python-klassen, 
waarbij de integriteit wordt afgedwongen via \textit{foreign keys} en een \textit{association table} voor de veel-op-veel relatie tussen concerten en artiesten. 
Dit zorgt ervoor dat data-consistentie behouden blijft wanneer bijvoorbeeld een artiest aan meerdere concerten wordt gekoppeld.

\begin{lstlisting}[language=Python,label={Concert Model},caption={Voorbeeld van een SQLAlchemy datamodel voor een concert.}]
def utcnow():
    return datetime.now(UTC).replace(tzinfo=None)

class Concert(Base):
    __tablename__ = "concerts"
    id = Column(Integer, primary_key=True, index=True)
    title = Column(String, nullable=False)
    date = Column(DateTime, nullable=False) # Afgeleid van Graves (2003): startTime
    price = Column(Float, nullable=True) # Afgeleid van Graves (2003): ticketPrice
    image_url = Column(String, nullable=True)
    url = Column(String, index=True, nullable=True) 
    venue_id = Column(Integer, ForeignKey("venues.id"), nullable=False) # Afgeleid van Graves (2003): Venue

    content_hash = Column(String, index=True, nullable=True)
    status = Column(String, default="active", nullable=False)
    last_scraped_at = Column(DateTime, default=utcnow, onupdate=utcnow)

    venue = relationship("Venue", back_populates="concerts")
    performers = relationship("Performer", secondary=concert_performers, back_populates="concerts") # Afgeleid van Graves (2003): Performer

    performer_ids = association_proxy("performers", "id")
\end{lstlisting}

\subsection{Datavalidatie en Schema's (Pydantic)}

Naast de database-modellen wordt Pydantic gebruikt voor datavalidatie en conversie bij API-verzoeken. 
De schema's bevatten de juiste regels om de integriteit van de data te bewaken tijdens de omzetting naar objecten. 
Dit is essentieel om ervoor te zorgen dat corrupte data de processen niet verstoren en dat de data consistent blijft tijdens de verwerking.

\begin{lstlisting}[language=Python,label={Concert Schema},caption={Voorbeeld van een Pydantic schema voor een concert.}]
class ConcertBase(BaseModel):
    title: str
    date: datetime
    price: Optional[float] = None
    image_url: Optional[str] = None
    url: Optional[str] = None
    content_hash: Optional[str] = None
    status: str = "active"
    last_scraped_at: Optional[datetime] = None

class Concert(ConcertBase):
    id: int
    created_at: datetime
    updated_at: datetime
    venue_id: int
    venue: Venue
    performers: List[Performer] = []

    model_config = ConfigDict(from_attributes=True)
\end{lstlisting}

\section{Extractie-mechanisme}

Om de architecturale schaalbaarheid te waarborgen, is er een configuratie-gestuurde aanpak gebruikt voor het extraheren van data van de verschillende venue-websites.
Zo kan de scraping-logica uniform worden toegepast op de verschillende HTML-structuren van de venue pagina's, origineel was dit met een .yaml configuratie, 
maar om de gebruiksvriendelijkheid te verhogen is dit verplaatst naar een eigenschap in de databank. Hierdoor is het mogelijk om met een simpel admin dashboard venues toe te voegen en te beheren zonder dat er code/configuratie manueel aangepast moet worden.

De extractie-logica bestaat hierdoor uit 2 delen:

\subsection{Scraper Engine}

De kern van de extractie vindt plaats in de \texttt{ScraperEngine}, die gebruikmaakt van de \texttt{Scrapling}-bibliotheek. Om de schaalbaarheid van de applicatie te garanderen, bevat deze engine geen hardgecodeerde regels voor specifieke zalen. In plaats daarvan fungeert het als een generieke motor die wordt aangedreven door een extern configuratie-object (\texttt{VenueScraperConfig}). 

Daarnaast is de engine ontworpen met ethisch scrapen in het achterhoofd: bij de initialisatie van de \texttt{Fetcher} wordt de naleving van \texttt{robots.txt} strikt afgedwongen. 

Listing \ref{lst:scraper-init} toont de opzet en de primaire iteratielus van de engine, waarbij te zien is hoe de inkomende configuratie de extractie aanstuurt.

\begin{lstlisting}[language=Python, label={lst:scraper-init}, caption={Initialisatie en configuratie-gestuurde extractie in de Scraper Engine.}]
class ScraperEngine:
    def __init__(self):
        # Ethisch scrapen: automatische naleving van robots.txt
        self.fetcher = Fetcher(auto_match=True, robots_txt_obey=True)

    def scrape_venue(self, config: VenueScraperConfig) -> List[Dict[str, Any]]:
        url = str(config.start_url)
        try:
            response = self.fetcher.get(url)
            # Dynamische selectie van elementen op basis van de configuratie
            cards = response.css(config.selectors.card)
        except Exception as e:
            logger.error(f"Extractie gefaald voor {config.venue_name}: {e}")
            return []
        
        scraped_data = []
        for card in cards:
            # Voorbeeld van dynamische extractie:
            title_sel = card.css(config.selectors.title)
            if not title_sel:
                continue
            title = " ".join(title_sel[0].get_all_text().split()).strip()
            
            # ... (extractie van datum, url, image en prijs via fallbacks)
\end{lstlisting}

Naast de ruwe data-extractie is de \texttt{ScraperEngine} ook verantwoordelijk voor de voorbereiding op change detection. Omdat websites geen unieke ID's meesturen voor hun evenementen, berekent de engine een cryptografische vingerafdruk (hash) op basis van de inhoudelijke velden. 

Zoals te zien in Listing \ref{lst:scraper-hash}, wordt deze \texttt{content\_hash} samen met de geëxtraheerde data gevalideerd via een Pydantic-model (\texttt{ScrapedItem}) voordat deze de engine verlaat. Dit garandeert dat onvolledige of corrupte iteraties de database nooit bereiken.

\begin{lstlisting}[language=Python, label={lst:scraper-hash}, caption={Generatie van de content hash en strikte datavalidatie.}]
                # Berekening van een unieke vingerafdruk voor change detection
                content_str = f"{title}|{dt.isoformat()}|{sorted(performers)}|{price}|{image_url}"
                content_hash = hashlib.md5(content_str.encode()).hexdigest()

                # Gebruik Pydantic voor finale validatie en opschoning
                item = ScrapedItem(
                    title=title,
                    date=dt,
                    url=full_url,
                    image_url=image_url,
                    price=price,
                    performers=performers,
                    venue_name=config.venue_name,
                    content_hash=content_hash
                )
                scraped_data.append(item.model_dump())
            except Exception as e:
                logger.error(f"Fout bij het parsen van item in {config.venue_name}: {e}")
                
        return scraped_data
\end{lstlisting}

\subsection{Scraping configuratie}

Elke scraper configuratie bevat subschema's die instructies bevatten over hoe de data van die specifieke venue moet worden geëxtraheerd.

\begin{lstlisting}[language=Python,label={Venue Scraper Configuratie},caption={Voorbeeld van een scraper configuratieschema voor een venue.}]
class VenueScraperConfig(BaseModel):
    venue_name: str 
    start_url: HttpUrl                      # Base URL van de venue website
    selectors: SelectorConfig               # CSS-selectors voor het extraheren van concertgegevens
    date_parsing: DateParsingConfig         # Instructies voor het parsen van datums
    performer_strategy: PerformerStrategy   # Instructies voor het extraheren van performer-informatie
\end{lstlisting}

\subsection{Omzeilen van Barrières}

Een essentieel onderdeel van de scraping-engine is het vermogen om de geavanceerde barrières van moderne concertwebsites te omzeilen. Hiervoor maakt de PoC gebruik van de \textit{stealth} mogelijkheden van Scrapling, in combinatie met de \textit{browserforge} library. De implementatie focust op drie niveaus:

\begin{itemize}
    \item \textbf{Browser Impersonation}: Door gebruik te maken van geavanceerde fetchers worden browserspecifieke headers en JavaScript-eigenschappen automatisch gesynchroniseerd met die van een legitieme browser. Dit voorkomt dat de website het gebruik van een \textit{headless browser} detecteert.
    \item \textbf{TLS Fingerprinting}: De engine simuleert de cryptografische patronen van echte browsers tijdens de SSL/TLS-handshake. Hierdoor wordt voorkomen dat de infra-laag van de website het verzoek blokkeert op basis van een vingerafdruk die typisch is voor Python-scripts.
    \item \textbf{Behavioral Simulation}: Het systeem hanteert dynamische extractiepatronen om te voorkomen dat geautomatiseerde frequenties leiden tot een IP-blokkade, terwijl het tegelijkertijd de ethische richtlijnen van de doelwebsite respecteert.
\end{itemize}

Zoals te zien is in Listing \ref{lst:stealth-init}, worden deze complexe mechanismen in de architectuur geabstraheerd tijdens de initialisatie van de \texttt{ScraperEngine}. Door de parameter \texttt{auto\_match=True} mee te geven, activeert het framework autonoom de browser-impersonatie en TLS-spoofing voor elk uitgaand verzoek.

\begin{lstlisting}[language=Python, label={lst:stealth-init}, caption={Initialisatie van de ScraperEngine met ingebouwde stealth- en compliance-mechanismen.}]
from scrapling import Fetcher

class ScraperEngine:
    def __init__(self):
        # Initialisatie van de Fetcher
        # auto_match=True activeert Browserforge (header impersonation)
        # en curl-cffi (TLS fingerprint spoofing) om bot-detectie te omzeilen
        self.fetcher = Fetcher(
            auto_match=True, 
            robots_txt_obey=True # Behoudt de ethische compliance
        )
\end{lstlisting}

\subsection{Strategische Selectie van de fetcher}
\label{sec:extractie-motor}

Het \texttt{Scrapling}-framework biedt een modulair arsenaal aan 
\textit{fetchers}, elk ontworpen voor specifieke 
weerstandsniveaus en performantie-eisen. Binnen de architectuur van 
deze \gls{PoC} is de extractielaag zodanig abstract dat het 
wisselen tussen deze motoren beperkt tot het aanpassen van een import en klasse naam
, zonder dat de parsing-logica of de 
CSS-selectors herschreven moeten worden. Dit draagt fundamenteel bij 
aan de architecturale schaalbaarheid van het systeem.

Voor de huidige implementatie is een strategische afweging gemaakt 
tussen de beschikbare motoren op basis van snelheid, resourcegebruik 
en capaciteiten om anti-bot maatregelen te omzeilen. De volgende fetchers zijn overwogen\autocite{ScraplingFetchers}:

\begin{itemize}
    \item \textbf{\texttt{Fetcher} (De Huidige Keuze):} 
    Deze fetcher opereert puur op de netwerklaag. Door gebruik te maken 
    van \texttt{curl-cffi} en de \texttt{auto\_match=True} parameter, 
    simuleert deze fetcher de TLS-vingerafdrukken van legitieme 
    browsers. Het resourcegebruik is extreem laag en de 
    verwerkingssnelheid is optimaal. Voor locaties die \gls{SSR} gebruiken, is dit de meest efficiënte oplossing.
    
    \item \textbf{\texttt{AsyncFetcher} (Voor Toekomstige Schaalbaarheid):} 
    Dit is de asynchrone variant van de standaard fetcher, ontworpen 
    om meerdere pagina's parallel op te halen via \texttt{asyncio}. 
    Omdat de scope van deze \gls{PoC} zich beperkt tot de validatie 
    op een handvol locaties, is de prestatiewinst in dit geval verwaarloosbaar. 
    In een productie-omgeving met honderden venues 
    zou deze fetcher de totale looptijd verlagen.
    
    \item \textbf{\texttt{DynamicFetcher} \& \texttt{StealthyFetcher} (Applicatielaag):} 
    Deze fetchers starten een volwaardige \textit{headless browser} op. 
    De \texttt{DynamicFetcher} dient om JavaScript-gedreven 
    \gls{SPA}s te renderen. De \texttt{StealthyFetcher} 
    voegt daar nog geavanceerde browser-impersonatie en canvas-spoofing aan toe. 
    Omdat het initiëren van browser-instanties veel geheugen vereist, 
    worden deze fetchers in deze \gls{PoC} niet gebruikt. Deze fetchers zijn bedoeld voor websites met agressieve JavaScript-gebaseerde botdetectie. 
    En de gekozen venues maken hier geen gebruik van, waardoor het extra gebruik niet gerechtvaardigd is.
\end{itemize}

Tabel \ref{tab:fetcher-vergelijking} biedt een overzichtelijk 
vergelijkingskader van deze afweging.

\begin{table}[h]
\centering
\caption{Vergelijkende analyse van Scrapling extractiemotoren}
\label{tab:fetcher-vergelijking}
\begin{tabular}{p{3cm} p{2.5cm} p{4cm} p{2.5cm}}
\hline
\textbf{Fetcher} & \textbf{Laag} & \textbf{Defensieve Kracht} & \textbf{Resourcekost} \\
\hline
\texttt{Fetcher} & Netwerk & Omzeilt TLS/JA3 vingerafdrukken & Zeer laag \\
\texttt{AsyncFetcher} & Netwerk & Parallelle verwerking & Laag \\
\texttt{DynamicFetcher} & Applicatie & Voert JavaScript uit & Hoog \\
\texttt{StealthyFetcher} & Applicatie & Omzeilt JS-fingerprinting & Zeer hoog \\
\hline
\end{tabular}
\end{table}

\section{\textit{Change Detection} logica}

De kern van de \textit{change detection} logica bevindt zich in een asynchrone Celery-taak die de gescrapete data vergelijkt met de huidige staat in de database . 
Door gebruik te maken van een \textit{content hash} kan het systeem efficiënt bepalen of de informatie van een concert gewijzigd is, zonder dat elke kolom manueel vergeleken moet worden. 
Voor een gedetailleerd overzicht van de implementatie van deze logica wordt verwezen naar de \nameref{ch:bijlage} (zie Listing~\ref{lst:change-detection}).

\section{Notificatiesysteem}

Het notificatiesysteem in deze \gls{PoC} heeft volgende eigenschappen:
\begin{itemize}
    \item \textbf{Notificatie-interval}: Omdat Vlaamse concerten vaak aangekondigd worden voor de ticketverkoop en deze applicatie niet ontworpen is als een ticketdistributieplatform, 
    is er gekozen voor een notificatie-interval van 24 uur om notificatievermoeidheid te voorkomen. De notificaties worden dus gebundeld en dagelijks verstuurd, in plaats van onmiddelijk bij elke detectie.
    \item \textbf{Medium}: In deze \gls{PoC} worden de notificaties verstuurd via e-mail, echter is het systeem modulair opgezet dus is het mogelijk om efficiënt andere notificatiemedia toe te voegen.
    \item \textbf{Inhoud}: De notificaties bevatten de naam van het concert, de datum, de venue en een link naar de originele concertpagina zodat het geïnteresseerde publiek toch naar de bronorganisatie geleid wordt. Dit is ook een bewuste keuze zodat de data van de bronorganisatie zoveel mogelijk behouden blijft en de impact op de bronorganisatie minimaal is zoals besproken in \nameref{subsec:ethiek-legaliteit}.
\end{itemize}

Het systeem is modulair opgezet, het bestaat uit 4 delen die samen een volwaardig notificatiesysteem vormen:

\begin{itemize}
    \item \textbf{Trigger}: In de Celery-taak die de \textit{change detection} uitvoert, wordt de notificatie-queue aangeroepen.
                            Dit gebeurt in \texttt{process\_scraped\_items} (Listing~\ref{lst:change-detection}) in de regels waar \texttt{queue\_notifications\_for\_concert} wordt aangeroepen, wat zowel bij updates als bij nieuwe concerten is.
    \item \textbf{Matcher}: De matcher is verantwoordelijk voor het toevoegen van notificaties aan de queue op basis van de gebruikersvoorkeuren. 
                            Deze functie \newline \texttt{queue\_notifications\_for\_concert} (Listing~\ref{lst:notification-queueing}) controleert of er een gebruikersvoorkeur matcht met de venue of artiest van het concert.
                            Indien dit zo is, wordt er gekeken of er al een notificatie hiervoor bestaat, als dat niet het geval is wordt er een nieuwe notificatie toegevoegd aan de queue. 
    \item \textbf{Queue}: De notificatiequeue is in deze \gls{PoC} geïmplemented als een database-tabel waarin de notificaties worden opgeslagen met een status aangezien deze gebundeld worden. Het gebruikte model voor deze queue staat in Listing~\ref{lst:notification-model}.
                          De Celery-worker die verantwoordelijk is voor het versturen van notificaties, haalt de items uit deze queue en verwerkt ze.
    \item \textbf{Mailer}: De mailer is verantwoordelijk voor het dagelijks versturen van de notificaties. Dit is geïmplementeerd als een Celery-task (Listing~\ref{lst:daily-digest-processor}) die elke 24 uur gepland wordt om alle nodige notificaties te versturen via de helper functie \texttt{send\_daily\_digest\_email} (Listing~\ref{lst:daily-digest-mailer}).
                            De Celery-task zorgt ervoor dat de notificaties worden gebundeld per gebruiker en dat de status van de notificaties in de queue correct wordt bijgewerkt na verzending. De helper functie is verantwoordelijk voor het opstellen en versturen van de e-mail, deze zorgt zowel voor een plain text als een HTML-versie van de e-mail zodat deze altijd kan worden weergegeven ondanks het gebruik van andere e-mailclients.
\end{itemize}

\section{API-implementatie en Beveiliging}

De \gls{API} is asynchroon opgezet met FastAPI, zodat de server responsief blijft en deze ook de scraper en Celery-taken gelijktijdig op de achtergrond kunnen draaien, zonder de gebruikerservaring te beïnvloeden.

\subsection{Modulaire Endpoints}

Elke entiteit heeft zijn eigen set van endpoints voor CRUD-operaties in de /api/endpoints folder. Deze isolering houdt de codebase overzichtelijk en onderhoudbaar, het scheidt de verschillende domeinen en maakt het eenvoudiger om eventuele uitbreidingen of wijzigingen door te voeren.
De uiteindelijke samenhang van deze endpoints wordt gevormd in de /api folder, waar de router als volgt wordt ingesteld:

\begin{lstlisting}[language=Python,label={API Router},caption={Voorbeeld van het samenstellen van de API-router in FastAPI.}]
    api_router = APIRouter()

    api_router.include_router(venues.router, prefix="/venues", tags=["venues"])
    api_router.include_router(concerts.router, prefix="/concerts", tags=["concerts"])
    api_router.include_router(performers.router, prefix="/performers", tags=["performers"])
    api_router.include_router(users.router, prefix="/users", tags=["users"])
    api_router.include_router(subscriptions.router, prefix="/subscriptions", tags=["subscriptions"])
\end{lstlisting}

Overigens wordt in main.py de router toegevoegd aan de FastAPI-applicatie, waardoor deze beschikbaar wordt voor communicatie met de frontend:

\begin{lstlisting}[language=Python,label={FastAPI Integratie App},caption={Voorbeeld van het instellen van de FastAPI-applicatie.}]
    from fastapi import FastAPI

    app = FastAPI(title="Concert Monitor API")

    # Include all endpoints
    app.include_router(api_router, prefix="/api/v1")
\end{lstlisting}

\subsection{Validatie}

Elk API-endpoint maakt gebruik van Pydantic-schema's voor datavalidatie en type-veiligheid. Deze schema's bevatten validatie voor elk veld, zoals het controleren van datums, e-mailformaten en verplichte velden. 
Zoals hieronder zichtbaar wordt de data telkens in de schema's gevalideerd voordat deze wordt verwerkt, waardoor de integriteit van de data wordt gewaarborgd en fouten vroegtijdig worden opgevangen.

\begin{lstlisting}[language=Python,label={API Endpoint Validatie},caption={Voorbeeld van een API-endpoint met Pydantic-validatie.}]
def get_active_concerts(db: Session, skip: int = 0, limit: int = 100, search: str = None):
    """
    Retrieve concerts where date >= now() AND status == 'active'.
    Includes venue and performers using joinedload.
    """
    query = select(Concert).options(joinedload(Concert.venue), joinedload(Concert.performers))
    
    filters = [
        Concert.date >= datetime.now(),
        Concert.status == "active"
    ]
    
    if search:
        # Search in title or venue name
        query = query.join(Concert.venue)
        filters.append(
            sa.or_(
                Concert.title.ilike(f"%{search}%"),
                Venue.name.ilike(f"%{search}%")
            )
        )
        
    return db.execute(
        query.where(and_(*filters))
        .offset(skip)
        .limit(limit)
    ).scalars().unique().all()
\end{lstlisting}

Overigens wordt er ook gebruik gemaakt van de functie \texttt{get\_db} in de API-endpoints, 
deze dependency injection waarborgt het lifecycle-beheer van de database-sessies.

\begin{lstlisting}[language=Python,label={Database Dependency Injection},caption={Voorbeeld van database dependency injection in FastAPI.}]
    def get_db():
    db = SessionLocal()
    try:
        yield db
    finally:
        db.close()
\end{lstlisting}

\subsection{Beveiliging}

Secrets zoals database SMTP credentials worden beheerd via omgevingsvariabelen en geladen met behulp van de Python-dotenv bibliotheek. 
Deze aanpak voorkomt dat gevoelige informatie in de codebase terechtkomt en maakt het gemakkelijker om verschillende configuraties te beheren voor ontwikkel-, test- en productieomgevingen.

Verder wordt er strikt gecontroleerd welke domeinen toegang hebben tot de \gls{API}, door middel van de CORS-middleware in FastAPI.

\begin{lstlisting}[language=Python,label={CORS Middleware},caption={Voorbeeld van het instellen van CORS-middleware in FastAPI.}]
    from fastapi.middleware.cors import CORSMiddleware

    app.add_middleware(
        CORSMiddleware,
        allow_origins=["*"],    # Dit hoort de toegestane domeinen te bevatten in een productieomgeving
        allow_credentials=True,
        allow_methods=["*"],
        allow_headers=["*"],
    )
\end{lstlisting}

